\chapter{Entities}

Just as there are many game engines, there are just as many definitions and nomenclature
for what to call the pieces that build up the game world.
Game worlds usually have many different types of objects,
but the most common distinction that can be made is between static and dynamic objects.
Static objects form up the unchangeable parts of the world, such as the ground, walls, and other objects that are not meant to be interacted with.
The other category is dynamic objects, these are the objects that can be interacted with by the player or other dynamic objects.
Despite the lack of a standardized terminology,
understanding the unique approaches taken by existing engines is important when developing a new engine.
Unity refers to its in-game objects as "GameObjects",
Unreal Engine 4/5 employs the term "Actors",
and the Source Engine designates its elements as "Entities".
While they use different names and have different capabilities,
the fundamental concepts behind them are the same.

Unity has GameObjects that build up everything in the game world.
These form a spatial hierarchy [1] that allows us to group and move game objects together,
this can be used for a player where we might need to move multiple models together.

<Image Hierarchy>

The user then can attach components to these game objects, which than define what it is capable of.
For example, the user can attach a Mesh and mesh renderer to create something that will show u in the world.
On to of the built-in components, the user can define their own components in c\# scripting by extending Mono behaviour

In Unreal Engine 4/5, Actors serve as the foundation for all interactive elements within a game world.
Similar to Unity's game objects, Actors form a hierarchical structure,
enabling the organization and manipulation of multiple objects in tandem,
such as character models or environmental assets.

<Image Outliner>

Within each Actor, Components can be added, which determine the object's functionality and behaviour.
These Components themselves can be organized in a hierarchical fashion,
allowing for a more structured and efficient setup.
For example, a Static Mesh Component can be added to an Actor, enabling it to be visually represented within the game world.
Users can also create custom Components using C++ or Blueprints,
Unreal Engine's visual scripting system, by extending the base Actor Component class.
This offers a powerful and flexible way to tailor functionality to specific needs while maintaining a well-organized Component hierarchy within each Actor.

Source engine uses a fix set of entities that are compiled with the engine.

\subsection{Shadow Engine}
In shadow engine we use a more complex setup that borrows form UE4/5.
In this solution we imagine the game world as a graph, namely a tree.
Everything in this tree is a Node, and these can be further specialized according to what role they fulfil.

There are 2 types of main Nodes: Actros and Entites, the difference is that Actors have names, and both can have public and internal hierarchy.
The internal hieracrhy is similar to what UE4/5 offers

The start of the graph is a RootNode that has all the loaded scenes as it's children.

entity graph
NodeBase
Node
Entity
Scene
	 